\newpage\thispagestyle{hl}
\hlbanner{Part II $\cdot$ Full-order physics}{3}{Linear continua and structural elements}
\hlsection{The question}
\begin{hlquote}A thin steel strip is meshed with a standard plate element, loaded at the tip and solved. The answer comes out two hundred times too stiff, and no setting is wrong: the integration rule that looks the most careful is the cause.\end{hlquote}
\vspace{8pt}\begin{hlsum}\textbf{The cause is the integration rule that looks most careful, and selective reduced integration removes it.} On the thin strip ($h/L=0.01$) the full-integration plate gives 0.005 of the beam-theory deflection on 4 elements and 0.602 on 64; with selective reduced integration it gives 1.059 and 1.078. The chapter goes from the weak form through shape functions and quadrature to bar, beam, plate, shell and solid elements, and compares consistent and lumped mass. The patch test does not catch this kind of locking, since every element passes it.\end{hlsum}
\hlsection{Key concepts}
\begin{description}\setlength\itemsep{5pt}\setlength\parsep{0pt}\renewcommand{\descriptionlabel}[1]{\hspace{\labelsep}\sffamily\bfseries\color{hlnavy}#1}
\item[Weak form] equilibrium tested against every admissible virtual displacement
\item[Shape functions, $\bm B$, $\bm D$] interpolation, strain from nodal values, stress from strain
\item[Gauss quadrature] the integration rule is a modelling choice, not only a numerical one
\item[Locking] spurious stiffening: shear in thin plates, volume near incompressibility
\item[Selective reduced integration] full rule for bending, reduced rule for the term that locks
\item[Patch test] a constant-strain state must be reproduced on distorted meshes
\end{description}
\hlsection{After this chapter you can}
\begin{itemize}\setlength\itemsep{3pt}
\item say how a weak form becomes a stiffness and a mass matrix
\item recognise locking or a wrong integration rule from a test with an independent reference
\item choose between consistent and lumped mass, knowing the frequency error of each
\end{itemize}
\newpage\thispagestyle{hl}
\noindent{\sffamily\small\color{hlgrey}Chapter 3 \quad Linear continua and structural elements}\par
\hlfig{\hlimg{ch03_bending}}{Pure bending of a plane-stress beam: tip deflection divided by the exact value against the number of elements. Linear elements converge from below; quadratic elements are exact.}
\hlsection{Headline result}
\begin{hlheadline}On a strip with $h/L=0.01$, the plate element with full integration gives 0.005 of the beam-theory deflection on 4 elements and still 0.602 on 64. With selective reduced integration it gives 1.059 and 1.078.\end{hlheadline}
\hlsection{From the code}
\begin{lstlisting}[style=hlpy]
s = FESystem(mesh, Quad4MindlinPlate())
s.assemble_stiffness((Db, Ds), integration=rule)   # rule = "full" or "sri"
# 4 elements:  full 0.005   selective-reduced 1.059
# 64 elements: full 0.602   selective-reduced 1.078
\end{lstlisting}
\hlsection{A stated limit}
\begin{hllimit}Locking is not removed by the patch test: every element passes it, and the plain 8-node hexahedron still locks.\end{hllimit}
